\chapter{Ecuaciones clásicas de la Física Matemática}
\begin{description}
    \item [Ecuación de ondas.] Versión en dimensión 1 en espacio:
        \begin{equation*}
            \frac{\partial^2 u}{\partial t^2} - c^2 \frac{\partial^2 u}{\partial x^2} = 0
        \end{equation*}
    \item [Ecuación del calor.] Versión en dimensión 1 en espacio:
        \begin{equation*}
            \frac{\partial u}{\partial t} - D\frac{\partial^2 u}{\partial x^2} = 0
        \end{equation*}
    \item [Ecuación de Laplace.] Versión en dimensión 2 en el espacio:
        \begin{equation*}
            \frac{\partial^2 u}{\partial x^2} + \frac{\partial^2 u}{\partial y^2} = 0
        \end{equation*}
\end{description}
Dado un campo vectorial $F:\mathbb{R}^d\to \mathbb{R}^d$, podemos escribir $F=(F_1,\ldots,F_d)$ con cada $F_i:\mathbb{R}^d\to \mathbb{R}$.

\begin{definicion}
    Para $F:\mathbb{R}^d\to \mathbb{R}^d$ definimos la divergencia de $F$ como la aplicación $div F:\mathbb{R}^d\to \mathbb{R}$ dada por:
    \begin{equation*}
        div F = \sum_{i=1}^{d} \frac{\partial F_i}{\partial x_i}
    \end{equation*}
    A veces se le denota por $\nabla\cdot F$, ya que formalmente tenemos:
    \begin{equation*}
        \left(\frac{\partial }{\partial x_1}, \ldots, \frac{\partial }{\partial x_d}\right) \cdot (F_1,\ldots,F_d)
    \end{equation*}
\end{definicion}

\noindent
Si pensamos que $F$ es un campo de fuerzas, $div F$ nos dice cómo varían los volúmenes en dicho campo de fuerzas.

\begin{ejemplo}
    Si $x\in \mathbb{R}^d$ se escribe como $x=(x_1,x_2,\ldots,x_d)$, podemos considerar $F(x)=x$ y tendríamos entonces que:
    \begin{equation*}
        div F = \frac{\partial x_1}{\partial x_1} + \frac{\partial x_2}{\partial x_2} + \ldots + \frac{\partial x_d}{\partial x_d} = d
    \end{equation*}
    Dada $f:\mathbb{R}^d\to \mathbb{R}$, podemos calcular $\nabla f:\mathbb{R}^d\to \mathbb{R}^d$ y:
    \begin{equation*}
        div(\nabla f) = \sum_{i=1}^{d} \frac{\partial }{\partial x_i}\left(\frac{\partial f}{\partial x_i}\right) = \sum_{i=1}^{d} \frac{\partial^2 f}{\partial x_i^2} = tr(D^2 f)
    \end{equation*}
    A esta cantidad se le suele denotar por $\triangle f$, o como $\nabla^2 f$, el ``Laplaciano de $f$''.
\end{ejemplo}

\noindent
En dimensiones arbitrarias tenemos:
\begin{description}
    \item [Ecuación del calor.] 
        \begin{equation*}
            \frac{\partial u}{\partial t} = D\triangle u
        \end{equation*}
    \item [Ecuación de ondas.] 
        \begin{equation*}
            \frac{\partial^2 u}{\partial t^2} = c^2 \triangle u
        \end{equation*}
    \item [Ecuación de Laplace.]
        \begin{equation*}
            \triangle u = 0
        \end{equation*}
        Ecuación de Poisson para $\rho$ dada:
        \begin{equation*}
            \triangle u = \rho(x)
        \end{equation*}
\end{description}
Aquí tenemos el convenio de que para $u = u(t,x_1,\ldots,x_d)$:
\begin{equation*}
    \triangle u = \sum_{i=1}^d \frac{\partial^2 u}{\partial x_i^2}
\end{equation*}
Es decir, sólo afecta a las variables espaciales.\\

\noindent
La ecuación de ondas y de Laplace son variacionales\footnote{Esto se debe a que en la ecuación de calor hay disipación de la energía.}, es decir, son la ecuación de Euler-Lagrange de determinado lagrangiano.\\

\noindent
Dada $u:\Omega\to \mathbb{R}$ con $\Omega\subset \mathbb{R}^d$ abierto, consideramos un funcional de la forma
\begin{equation*}
    \cc{F}(u) = \int_\Omega F(x,u,\nabla u)~dx \qquad (\text{donde\ } x=(x_1,\ldots,x_d))
\end{equation*}
La ecuación de Euler-Lagrange asociada es en $\Omega$:
\begin{equation}\label{eq:abreviar}
    F_u - \frac{d}{dx_1}(F_{p_1}) - \frac{d}{dx_2}(F_{p_2}) - \ldots - \frac{d}{dx_d}(F_{p_d})
\end{equation}
Donde hemos escrito:
\begin{equation*}
    F_u = \frac{\partial F}{\partial u}, \qquad F_{p_i} = \frac{\partial F}{\partial p_i} 
\end{equation*}
Con el convenio de que $F=F(x,u,p_1,\ldots,p_d)$, donde $p_i = \frac{\partial u}{\partial x_i}$. Observamos que $F:\Omega\times\mathbb{R}\times\mathbb{R}^d\to \mathbb{R}$.\\

\noindent
En vista de la definición de la divergencia, podemos resumir~\eqref{eq:abreviar} como:
\begin{equation*}
    F_u - div_x(F_p) = 0
\end{equation*}
Donde entendemos que $F_p = (F_{p_1}, \ldots, F_{p_d})$.

\noindent
Para $d=2$ variables esto funciona porque (variable $(x,y)$):
\begin{align*}
    \delta_\varphi \cc{F}(u) &= \int_\Omega F_u\varphi + F_{p_x}\frac{\partial \varphi}{\partial x} + F_{p_y} \frac{\partial \varphi}{\partial y}~dxdy = 0 \\
                             &= \int_\Omega F_u\varphi - \frac{d}{dx}(F_{p_x})\varphi - \frac{d}{dy}(F_{p_y})\varphi~dxdy = 0 \\
                             &= \int_\Omega(F_u - \frac{d}{dx}(F_{p_x}) - \frac{d}{dy}(F_{p_y}))\varphi~dxdy = 0
\end{align*}

\begin{ejemplo}
    Sea $S$ una superficie descrita mediante la gráfica de una función $u:\Omega\subset\mathbb{R}^2\to \mathbb{R}$. 
\end{ejemplo}
La ecuación de Euler-Lagrange es:
\begin{equation*}
    F_u - \frac{d}{dx}(F_{p_x}) - \frac{d}{dy}(F_{p_y}) = 0 \qquad \text{en\ } \Omega
\end{equation*}
Donde:
\begin{equation*}
    F_{p_x} = \frac{\partial F}{\partial p_x}, \quad F_{p_y} = \frac{\partial F}{\partial p_y}, \quad p_x = \frac{\partial u}{\partial x},\quad p_y = \frac{\partial u}{\partial y}
\end{equation*}
Tenemos:
\begin{equation*}
    F = F(x,y,u,p_x,p_y) = \sqrt{1+{(p_x)}^{2} + {(p_y)}^{2}}
\end{equation*}
Así:
\begin{equation*}
    F_{p_x} = \frac{p_x}{\sqrt{1+{(p_x)}^{2}+{(p_y)}^{2}}}, \quad F_{p_y} = \frac{p_y}{\sqrt{1+{(p_x)}^{2}+{(p_y)}^{2}}}
\end{equation*}
La ecuación queda:
\begin{equation*}
    -\frac{\partial }{\partial x}\left(\frac{\frac{\partial u}{\partial x}}{\sqrt{1+{\left(\frac{\partial u}{\partial x}\right)}^{2}+{\left(\frac{\partial u}{\partial y}\right)}^{2}}}\right) - \frac{\partial }{\partial y}\left(\frac{\frac{\partial u}{\partial y}}{\sqrt{1+{(\frac{\partial u}{\partial x})}^{2}+{\left(\frac{\partial u}{\partial y}\right)}^{2}}}\right) = 0
\end{equation*}
Es una ``ecuación de las superficies mínimas''. Podemos abreviarla como:
\begin{equation*}
    div\left(\frac{\nabla u}{\sqrt{1+|\nabla u|^2}}\right) = 0
\end{equation*}
